% Source: https://artofproblemsolving.com/community/c3h1660029_trigonometry
% Stable identity: aops-topic-1660029-q01
Let $s=\sin x$ and $c=\csc x=\dfrac{1}{s}$. For $0^\circ<x<90^\circ$ we have $0<s<1<c$.

\textbf{First terms.} For a base $s$ with $0<s<1$, the power $s^t$ decreases as $t$ increases. Since $s<c$, we get $s^s>s^c$.

\textbf{Second terms.} For a base $c>1$, the power $c^t$ increases with $t$. Since $c>s$, we get $c^c>c^s$.

Adding,
\[
L-M=\left(s^s-s^c\right)+\left(c^c-c^s\right)>0,
\]
so $\boxed{L>M}$ for every $x$ in the interval. (At $x=90^\circ$ we would have $s=c=1$ and $L=M=2$, which is why $90^\circ$ is excluded.)
